% ============================================================
% First-Shell-to-First-Shell Edge-Direction Perpendicularity
%   on the 600-Cell Edge Graph at Vertex-Aligned Reading C
% Publication-Grade Hardening of Patch 0544 §4.1 Theorem 4.1.1
% Conscious Point Physics Flagship Paper Series
%   (F.1 Dynamical Substrate Law sub-question hardened-theorem artifact;
%    second of three Layer 3 building-block hardenings)
% Patch 0551 (Session 142, 23 May 2026)
% ============================================================
% Version 1.1 --- 17 August 2026 (Patch 3214):
%   added a How-we-review methodology note stating plainly that AI panel
%   review is a self-applied verification method and not journal peer
%   review; twelve chirality papers retitled so the descriptive title
%   leads and the identifier is demoted to a subtitle. No claim,
%   equation, number or section changed.
% Version 1.0 --- 17 August 2026 (Patch 3213):
%   extended the generated identifier appendix to all five identifier
%   families (THEO-, FI-, PRED-, CONJ-, PH- as well as OPEN-), and
%   replaced review-convergence scores with AI review panel wording
%   where present. No claim, equation, number or section changed.
%   This is the FIRST version stamp assigned to this paper; 1.0 denotes
%   the first stamped revision, NOT a claim of v1.0 shipped grade.

\documentclass[11pt,letterpaper]{article}
\usepackage[utf8]{inputenc}
\usepackage[margin=1in]{geometry}
\usepackage[T1]{fontenc}
\usepackage{amsmath,amssymb,amsfonts,amsthm}
\usepackage{xcolor}
\usepackage{hyperref}
\usepackage{enumitem}
\usepackage{parskip}
\usepackage{mdframed}

\hypersetup{
    colorlinks=true,
    linkcolor=blue,
    citecolor=blue,
    urlcolor=blue
}

\newtheorem{theorem}{Theorem}[section]
\newtheorem{proposition}[theorem]{Proposition}
\newtheorem{lemma}[theorem]{Lemma}
\newtheorem{corollary}[theorem]{Corollary}
\newtheorem{definition}[theorem]{Definition}
\newtheorem{remark}[theorem]{Remark}

\newcommand{\nhat}{\hat{n}}
\newcommand{\vhost}{v_{\text{host}}}
\newcommand{\Gsixcell}{\Gamma_{600}}
\newcommand{\Reals}{\mathbb{R}}
\newcommand{\Hperp}{H_{\vhost}}

\title{First-Shell-to-First-Shell Edge-Direction Perpendicularity \\
in the 600-Cell at Vertex-Aligned Reading C \\
\large Publication-Grade Hardening of Patch~0544~\S4.1 Theorem~4.1.1 (formerly Identity G3)}

\author{Thomas Lee Abshier, ND, and Claude Opus 4.7\\Hyperphysics Institute --- F.1 Dynamical Substrate Law trajectory}

\date{23 May 2026 (Session 142, CPP Patch 0551)\\[2pt]
{\small Version~1.0 --- 17 August 2026 (Patch 3213: extended the generated identifier appendix to all five identifier families (THEO-, FI-, PRED-, CONJ-, PH- as well as OPEN-), and replaced review-convergence scores with AI review panel wording where present. No claim, equation, number or section changed.)}\\[2pt]
{\small Version~1.1 --- 17 August 2026 (Patch 3214: added a How-we-review methodology note stating plainly that AI panel review is a self-applied verification method and not journal peer review; twelve chirality papers retitled so the descriptive title leads and the identifier is demoted to a subtitle. No claim, equation, number or section changed.)}}

\begin{document}
\maketitle

\begin{abstract}
We give a publication-grade hardening of the first-shell-to-first-shell edge-direction perpendicularity identity for the 600-cell edge graph at vertex-aligned Reading~C, registered in the CPP F.1 sub-question trajectory at Patch~0544 \S4.1 as Theorem~4.1.1 (formerly named Identity~G3) at sketch-document Layer~3 rigor. The hardened statement asserts that, under (i)~vertex-aligned Reading~C with $\nhat = \vhost$, (ii)~the 600-cell unit-vertex normalisation $|v|=1$, and (iii)~the first-shell inner-product primitive G1 ($v_i\cdot\vhost = \phi/2$ uniformly across the 12 first-shell neighbours of $\vhost$) inherited from Patch~0541 \S3, every unit edge-direction $\hat{e}_{ij}$ between distinct first-shell vertices $v_i \neq v_j$ of $\vhost$ satisfies $\hat{e}_{ij}\cdot\nhat = 0$ exactly. The proof is a two-line algebraic computation supported by a tangent-hyperplane geometric lemma; publication-grade rigorisation consists in explicit hypothesis tracking, isolation of the geometric content as Lemma~\ref{lem:tangent_hyperplane}, and a five-class exclusion enumeration that registers G1's own sketch-document Layer~3 status (Patch~0541 \S3.1) as a separate publication-grade hardening candidate. This artifact is the second of three Layer~3 building-block hardenings needed for full publication-grade closure of the F.1 sub-question Layer~3 stack; the first (Patch~0550) hardened the perturbation-theory propagation rule (Patch~0544 \S3.3 Theorem~3.3.3 + Corollary~3.3.4), which cites the present theorem as a dependency.
\end{abstract}

\section{Introduction and statement of the main result}\label{sec:introduction}

This paper hardens the first-shell-to-first-shell edge-direction perpendicularity identity for the 600-cell edge graph at vertex-aligned Reading~C from sketch-document Layer~3 rigor (the two-line algebraic computation at Patch~0544 \S4.1) to publication-grade rigor with explicit hypothesis tracking and isolation of the underlying geometric content. The substantive theorem content is unchanged from Patch~0544 \S4.1; the contribution of this artifact is the rigorisation of the proof structure, the isolation of the tangent-hyperplane geometric content as Lemma~\ref{lem:tangent_hyperplane}, the formalisation of the genericity hypothesis (distinct first-shell vertex pairs), and an explicit exclusion enumeration.

\paragraph{Position in the F.1 Layer 3 hardening sequence.}
Patch~0550 hardened the perturbation-theory propagation rule (Patch~0544 \S3.3 Theorem~3.3.3 + Corollary~3.3.4) at \texttt{flagship\_papers/dynamical\_substrate\_law/hardened\_theorems/perturbation\_locality\_propagation.tex}; that hardening cites the present theorem (then at Patch~0544 \S4.1 sketch-document Layer~3 rigor) as exclusion class~E5 ``underlying graph-theoretic claims about $\Gsixcell$ taken as given.'' The present paper closes that hardening dependency for one of the two identities. The remaining dependency, Patch~0544 \S4.2 Theorem~4.2.1 (host-to-first-shell uniform projection, formerly Identity~I1), is the natural target of an anticipated Patch~0552.

\paragraph{Statement of the main result.}
Let $\Gsixcell$ denote the 600-cell edge graph (Patch~0550 \S2.1; \cite{coxeter1973}) with vertex set the 120 vertices of the 600-cell normalised to $|v|=1$ on $S^3 \subset \Reals^4$. Fix a host vertex $\vhost \in V(\Gsixcell)$ and write the 12 first-shell neighbours of $\vhost$ as $v_1, \ldots, v_{12}$.

\begin{theorem}[First-shell-to-first-shell perpendicularity, hardened]\label{thm:first_shell_perpendicularity}
Under the following hypotheses:
\begin{enumerate}[label=(H\arabic*)]
\item \textbf{Vertex-aligned Reading~C with framework chirality direction $\nhat = \vhost$}, as registered in Capotauro~v2.0 FI-C-RC-1 + FI-C-RC-2;
\item \textbf{600-cell unit-vertex normalisation}: every vertex $v \in V(\Gsixcell)$ satisfies $|v| = 1$ in the ambient $\Reals^4$;
\item \textbf{First-shell inner-product primitive (G1)}: for every first-shell vertex $v_i$ of $\vhost$ (each of the 12 nearest neighbours), $v_i \cdot \vhost = \phi/2$, where $\phi = (1+\sqrt{5})/2$ is the golden ratio. (G1 is derived at Patch~0541 \S3.1 at sketch-document Layer~3 rigor from 600-cell coordinate primitives + the icosahedral first-shell inner-product structure $\cos(36^\circ) = \phi/2$; the derivation is imported here.)
\end{enumerate}
For any pair of distinct first-shell vertices $v_i \neq v_j$ of $\vhost$, the unit edge-direction vector
\begin{equation}\label{eq:edge_direction_definition}
\hat{e}_{ij} \;=\; \frac{v_j - v_i}{|v_j - v_i|}
\end{equation}
satisfies
\begin{equation}\label{eq:first_shell_perpendicularity}
\hat{e}_{ij} \cdot \nhat \;=\; 0
\end{equation}
exactly.
\end{theorem}

The proof is organised in two steps: a geometric lemma isolating the tangent-hyperplane content (\S\ref{sec:tangent_hyperplane_lemma}), and the algebraic computation completing the proof (\S\ref{sec:perpendicularity_proof}). Corollary~\ref{cor:icosahedron_in_tangent} below records the tangent-hyperplane geometric content as a stand-alone consequence.

\begin{remark}[Distinctness hypothesis $v_i \neq v_j$]\label{rem:distinctness}
The hypothesis $v_i \neq v_j$ excludes the degenerate case in which~\eqref{eq:edge_direction_definition} would have a zero denominator. The 600-cell first-shell vertices are 12 geometrically distinct points; the hypothesis is a technical convention, not a structural restriction. Note also that Theorem~\ref{thm:first_shell_perpendicularity} as stated does NOT require the pair $(v_i, v_j)$ to be adjacent in $\Gsixcell$ (i.e.\ connected by a 600-cell edge); the result holds for all $\binom{12}{2} = 66$ distinct first-shell pairs, of which exactly 30 are 600-cell edges (forming the 30 edges of the first-shell icosahedron). The substrate-physics application at Patch~0544 \S4.3 concerns specifically the 30 first-shell-to-first-shell substrate edges.
\end{remark}

\section{Preliminaries and framework commitments}\label{sec:preliminaries}

\subsection{The 600-cell edge graph and the first-shell icosahedron}\label{sec:first_shell_setup}

The 600-cell is the regular 4-dimensional polytope with Schl{\"a}fli symbol $\{3,3,5\}$ \cite{coxeter1973}; we adopt the unit-radius normalisation $|v|=1$ for all 120 vertices on $S^3 \subset \Reals^4$, with edge length $|v_a - v_b| = 1/\phi$ for adjacent vertex pairs (Patch~0541 \S3.1 derivation). The 12 first-shell vertices of any host vertex $\vhost$ form a regular icosahedron in the 3D tangent hyperplane to the $S^3$ at $\vhost$ (Patch~0541 \S2.2 G2 primitive), with first-shell-to-first-shell adjacencies giving the 30 edges of that icosahedron.

\subsection{Vertex-aligned Reading C and the chirality direction}\label{sec:vertex_aligned_reading_c}

Under vertex-aligned Reading~C (Capotauro v2.0 FI-C-RC-1 + FI-C-RC-2), the framework chirality direction $\nhat$ is identified with the host vertex direction: $\nhat = \vhost$. Together with the unit-vertex normalisation (H2), this means $|\nhat| = 1$ and $\nhat$ is itself a vertex of the 600-cell.

\subsection{The primitive G1 inner-product identity}\label{sec:g1_primitive}

\begin{lemma}[G1, first-shell inner-product primitive; imported from Patch~0541 \S3.1]\label{lem:G1}
Under (H2) and the 600-cell geometric structure, for every first-shell vertex $v_i$ of $\vhost$ in the 600-cell,
\begin{equation}\label{eq:G1_identity}
v_i \cdot \vhost \;=\; \frac{\phi}{2}
\end{equation}
uniformly across $i \in \{1, \ldots, 12\}$.
\end{lemma}

\noindent\textbf{Source of Lemma~\ref{lem:G1}.} The identity~\eqref{eq:G1_identity} follows from the 600-cell first-shell-edge inner product equalling $\cos(36^\circ) = \phi/2$, where $36^\circ$ is the dihedral angle between adjacent vertex-directions at any 600-cell vertex; the derivation is given at Patch~0541 \S3.1 at sketch-document Layer~3 rigor from the 600-cell coordinate structure + the golden-ratio identity $\cos(36^\circ) = \phi/2$. The present paper imports Lemma~\ref{lem:G1} as a sketch-document Layer~3 result; promoting G1 to publication-grade rigor is registered as exclusion class~E1 of \S\ref{sec:scope_and_exclusions} and remains a candidate for a separate publication-grade hardening Patch.

\section{Proof of the perpendicularity theorem}\label{sec:proof}

\subsection{Tangent-hyperplane lemma}\label{sec:tangent_hyperplane_lemma}

\begin{lemma}[First-shell tangent hyperplane]\label{lem:tangent_hyperplane}
Under (H1)--(H3) of Theorem~\ref{thm:first_shell_perpendicularity}, every first-shell vertex $v_i$ of $\vhost$ ($i \in \{1, \ldots, 12\}$) lies in the 3-dimensional affine hyperplane
\begin{equation}\label{eq:tangent_hyperplane}
\Hperp \;=\; \bigl\{\, w \in \Reals^4 \,\big|\, w \cdot \nhat = \tfrac{\phi}{2}\,\bigr\} \;\subset\; \Reals^4.
\end{equation}
\end{lemma}

\begin{proof}
By Lemma~\ref{lem:G1} (G1), $v_i \cdot \vhost = \phi/2$ for each $i \in \{1, \ldots, 12\}$. Under hypothesis~(H1), $\nhat = \vhost$. Substituting gives $v_i \cdot \nhat = \phi/2$, so $v_i \in \Hperp$ by~\eqref{eq:tangent_hyperplane}. The argument is uniform in $i$. $\square$
\end{proof}

\begin{remark}[Geometric content of Lemma~\ref{lem:tangent_hyperplane}]\label{rem:tangent_hyperplane_geometric}
The affine hyperplane $\Hperp$ defined by~\eqref{eq:tangent_hyperplane} is the 3D affine subspace of $\Reals^4$ obtained from the 3D linear subspace $\nhat^\perp$ (orthogonal complement of $\nhat$) by translation along $\nhat$ by amount $\phi/2$. Equivalently, $\Hperp$ is the affine 3D tangent hyperplane to the sphere $S^3$ at the polar-cap height $\phi/2$ above the antipodal direction $-\nhat$. The 12 first-shell vertices, as a regular icosahedron centred on the $\nhat$-axis, lie on the 2-sphere $S^2 = S^3 \cap \Hperp$ of polar-cap radius $\sqrt{1 - (\phi/2)^2}$.
\end{remark}

\subsection{Perpendicularity computation}\label{sec:perpendicularity_proof}

\begin{proof}[Proof of Theorem~\ref{thm:first_shell_perpendicularity}]
Let $v_i, v_j$ be distinct first-shell vertices of $\vhost$. By Lemma~\ref{lem:tangent_hyperplane}, both $v_i$ and $v_j$ lie in $\Hperp$, so $v_i \cdot \nhat = v_j \cdot \nhat = \phi/2$.

Compute, using the linearity of the inner product:
\begin{align*}
\hat{e}_{ij} \cdot \nhat
\;=\;& \frac{(v_j - v_i) \cdot \nhat}{|v_j - v_i|} \\
\;=\;& \frac{(v_j \cdot \nhat) - (v_i \cdot \nhat)}{|v_j - v_i|} \\
\;=\;& \frac{(\phi/2) - (\phi/2)}{|v_j - v_i|} \\
\;=\;& \frac{0}{|v_j - v_i|}.
\end{align*}
By Remark~\ref{rem:distinctness}, $v_i \neq v_j$ gives $|v_j - v_i| > 0$, so the quotient is well-defined and equals $0$. Therefore $\hat{e}_{ij}\cdot\nhat = 0$, as required. $\square$
\end{proof}

\begin{corollary}[Icosahedral first shell lies in the tangent hyperplane]\label{cor:icosahedron_in_tangent}
The 12 first-shell vertices of $\vhost$ form a regular icosahedron contained in the 3D affine hyperplane $\Hperp$. Consequently the 30 first-shell-to-first-shell 600-cell edges (the edges of this icosahedron) all lie within $\Hperp$, and their unit edge-direction vectors lie in the 3D linear subspace $\nhat^\perp = \{w \in \Reals^4 : w\cdot\nhat = 0\}$.
\end{corollary}

\begin{proof}
The first-shell-icosahedron content follows from Patch~0541 \S2.2 primitive G2; Lemma~\ref{lem:tangent_hyperplane} places the icosahedron in $\Hperp$. For any pair of icosahedron vertices $v_i, v_j$, the edge $v_i v_j$ lies in $\Hperp$ because $\Hperp$ is affine (closed under convex combinations). The unit edge-direction $\hat{e}_{ij}$ is parallel to the difference $v_j - v_i$, which lies in the linear subspace $\nhat^\perp$ (the translation of $\Hperp$ by $-\phi\nhat/2$). By Theorem~\ref{thm:first_shell_perpendicularity}, $\hat{e}_{ij} \cdot \nhat = 0$, confirming $\hat{e}_{ij} \in \nhat^\perp$. $\square$
\end{proof}

\section{Scope, framework-locality, and explicit exclusions}\label{sec:scope_and_exclusions}

Theorem~\ref{thm:first_shell_perpendicularity} is, by design, scope-bounded: hypotheses~(H1)--(H3) collectively encode the framework commitments under which the perpendicularity identity holds. This section enumerates five classes of generalisations that lie explicitly outside the theorem's scope.

\begin{enumerate}[label=\textbf{(E\arabic*)}]
\item \textbf{Lemma~\ref{lem:G1} (G1) itself is imported at sketch-document Layer~3 rigor.} The first-shell inner-product primitive $v_i \cdot \vhost = \phi/2$ is derived at Patch~0541 \S3.1 from the 600-cell first-shell-edge dihedral angle $\cos(36^\circ) = \phi/2$ together with the unit-vertex normalisation. Promotion of Lemma~\ref{lem:G1} to publication-grade rigor is a candidate for a separate hardening Patch (the natural target being a publication-grade derivation of the 600-cell first-shell inner-product structure from the polytope's coordinate description).
\item \textbf{Layer 4 derivation of the 600-cell substrate from CPP axioms $A1$--$A11$.} The 600-cell polytope structure itself is taken here as given (H2 + Patch~0541 \S2.2 G1 + G2 primitives). A Layer~4 derivation of the 600-cell substrate from CPP axioms together with first-principles emergence considerations is the long-term programme target registered at Patch~0528 \S14.17 and carried forward through Patches~0539, 0549.
\item \textbf{Non-vertex-aligned Reading~C variants are NOT covered.} The theorem requires $\nhat = \vhost$ (hypothesis~(H1)). Alternative Reading~C placements --- centre-aligned Reading~C, face-aligned Reading~C, edge-aligned Reading~C --- would change which directions in $\Reals^4$ are identified as $\nhat$ and would generally invalidate Lemma~\ref{lem:tangent_hyperplane}. The Capotauro v2.0 FI-C-RC-2 specifies vertex-alignment; the present theorem is tied to that specification.
\item \textbf{Higher-shell perpendicularity claims are NOT addressed.} Theorem~\ref{thm:first_shell_perpendicularity} addresses only first-shell-to-first-shell edge directions. Second-shell and higher-shell perpendicularity properties are not covered; in fact, second-shell vertices lie in distinct affine hyperplanes parallel to $\Hperp$ at different polar-cap heights, so analogues of Theorem~\ref{thm:first_shell_perpendicularity} for second-shell-to-second-shell edges hold but with different normal vectors and would require separate proofs.
\item \textbf{The theorem is restricted to distinct first-shell vertex pairs.} The degenerate case $v_i = v_j$ is excluded by Remark~\ref{rem:distinctness}; while $\hat{e}_{ii}$ is undefined geometrically (there is no edge from a vertex to itself), this is a notational convention, not a structural restriction. The theorem applies to all $\binom{12}{2} = 66$ distinct first-shell pairs, irrespective of whether the pair is connected by a 600-cell edge. The 30 pairs that ARE 600-cell-edge-connected (forming the first-shell icosahedron) are the load-bearing case for the substrate-physics application at Patch~0544 \S4.3 and Patch~0550 Remark~3.3.2.
\end{enumerate}

\section{Cross-references and consequences}\label{sec:cross_references}

\subsection{Use in Patch 0550's perturbation-theory propagation rule}\label{sec:use_in_0550}

Patch~0550 \S3.3 Remark~3.3.2 (\texttt{hardened\_theorems/perturbation\_locality\_propagation.tex}) cited Patch~0544 \S4.1 Theorem~4.1.1 to justify the observation that perturbed steps along first-shell-to-first-shell edges in the 600-cell contribute zero to the $\mathcal{O}(\delta^n)$ coefficient of the substrate net DI-bit current. The present hardening closes that cross-reference dependency at publication-grade rigor: future content citing Patch~0550 Remark~3.3.2 may now cite the present Theorem~\ref{thm:first_shell_perpendicularity} in place of the unhardened Patch~0544 \S4.1 statement.

\subsection{Use in substrate-rate perturbation analysis}\label{sec:use_in_substrate_rate}

Applying Mechanism~A (Patch~0550 Definition~2.2.1) with Theorem~\ref{thm:first_shell_perpendicularity}: for any first-shell-to-first-shell edge $(v_i, v_j)$, the propagation rate
\begin{equation*}
r(\hat{e}_{ij}) \;=\; r_0\bigl(1 + \delta\,\hat{e}_{ij}\cdot\nhat\bigr) \;=\; r_0\bigl(1 + \delta\cdot 0\bigr) \;=\; r_0
\end{equation*}
is exactly $r_0$, independent of $\delta$ and uniform across all 30 first-shell-to-first-shell substrate edges. This is the ``zero first-order perturbation on first-shell-to-first-shell edges'' content of Patch~0544 \S4.3.

\subsection{Capotauro v2.0 §13.3 structural parallel}\label{sec:capotauro_parallel}

The geometric content of Theorem~\ref{thm:first_shell_perpendicularity} is identical to the Capotauro v2.0 \S13.3 spatial-sector ``structural proof'' parallel: under vertex-aligned Reading~C, the first-shell icosahedron lies in the 3D tangent hyperplane $\Hperp$ orthogonal to $\nhat$ (modulo the $\phi/2$ translation), and the edge-vectors within the icosahedron are tangent to this hyperplane and thus orthogonal to $\nhat$. The Capotauro v2.0 derivation operates in the spatial sector (perturbation amplitude $\epsilon$) and the present derivation operates in the temporal sector (perturbation amplitude $\delta$), but the underlying geometric content is identical and now hardened at publication-grade rigor in both sectors.

\section{Conclusion}\label{sec:conclusion}

Theorem~\ref{thm:first_shell_perpendicularity}, with hypotheses (H1)--(H3) and supporting Lemmas~\ref{lem:G1} (G1, imported) and~\ref{lem:tangent_hyperplane} (tangent hyperplane), gives the publication-grade hardening of Patch~0544 \S4.1 Theorem~4.1.1 (formerly Identity~G3). The scope qualifier of vertex-aligned Reading~C is built into hypothesis~(H1), the 600-cell unit-vertex normalisation into hypothesis~(H2), and the load-bearing G1 dependency into hypothesis~(H3). The five exclusion classes (E1)--(E5) make explicit what is NOT covered.

This artifact is the second of three Layer~3 building-block hardenings needed for full publication-grade closure of the F.1 sub-question Layer~3 stack. The remaining building-block hardening (Patch~0544 \S4.2 Theorem~4.2.1, host-to-first-shell uniform projection / Identity~I1) is the natural target of an anticipated Patch~0552; after Patch~0552 lands, the F.1 flagship paper assembly evaluation gate (LIVE since Patch~0549) becomes substantively engageable with three publication-grade theorems as building blocks: Patch~0550 (perturbation-theory propagation), the present Patch~0551 (first-shell-to-first-shell perpendicularity), and the anticipated Patch~0552 (host-to-first-shell uniform projection).

\bibliographystyle{plain}
\begin{thebibliography}{99}

\bibitem{coxeter1973}
H.~S.~M.~Coxeter,
\textit{Regular Polytopes},
Dover Publications, New York, 1973.

\bibitem{abshier_capotauro_chiral_mechanism_sketch}
T.~L.~Abshier,
``The Capotauro Mechanism v2.0: Chirality on the K3-Doublet from Substrate-Vacuum Broken-Symmetry Physics,''
CPP Flagship Paper Series, May~2026, OSF DOI \texttt{10.17605/OSF.IO/JXE8D}.

\end{thebibliography}

% BEGIN GENERATED IDENTIFIER APPENDIX -- do not edit by hand
\clearpage
\section*{Appendix: Programme identifiers used in this paper}
\addcontentsline{toc}{section}{Appendix: Programme identifiers}

\noindent This programme tracks its results, assumptions and unresolved questions under short identifiers, so that a claim and its dependencies can be cited precisely. Prefixes denote the kind of item: \texttt{THEO-} a proved theorem, \texttt{FI-} a foundational input the derivation assumes, \texttt{PRED-} a registered prediction, \texttt{CONJ-} a conjecture, \texttt{OPEN-} an unresolved question, \texttt{PH-} a problem history. Those appearing in this paper are glossed below; no external document is needed to read them.

\begin{description}
  \item[\texttt{FI-C-RC-1}] \hfill \\ The substrate primitive: a preferred 4D direction in the ambient space where the 600-cell substrate sits, sourcing all parity-sensitive observables.
  \item[\texttt{FI-C-RC-2}] \hfill \\ The vertex-aligned reading identifying that 4D direction with the host vertex, fixing the local structure as the icosahedral first shell.
\end{description}

\subsection*{How we review}

\noindent Results in this programme are checked by an \emph{AI review panel}: several large language models from different vendors are given the same paper and the same fixed protocol, and asked to find errors and raise objections independently. Objections are recorded and either answered in the text or registered as open problems under the identifiers glossed above.

\noindent This is a \emph{verification method the authors apply to their own work}. It is not journal peer review, it is not equivalent to it, and agreement among panel members is not evidence that a result is correct --- only that no member found a specific fault. Where individual models are named in the text, they are named for reproducibility of the method; model versions change, and a reader should treat the named models as a record of what was run rather than as an endorsement.

% END GENERATED IDENTIFIER APPENDIX

\end{document}
